$F \subset K$ -- расширение полей,  $\alpha \in K$

Следующие условия эквивалентны:

\begin{enumerate}
    \item $[F(\alpha):F]$ -- конечное
    \item $\exists f(x) \in F[x]$, т.ч. $f(\alpha) = 0$
\end{enumerate}

\Proof

\begin{itemize}
    \item В одну сторону довольно очевидно:
    
    Если  $[F(\alpha):F] = N < \infty$, то $1, \alpha, \ldots, \alpha^N$ ЛЗ над $F$. \newline Тогда $a_0 + a_1\alpha + \ldots a_N\alpha^N = 0, a_0, \ldots, a_N \in F$. \newline Тогда можно взять многочлен $f(x) = a_0 + a_1x + \ldots a_Nx^N$, который будет обнуляться для $x=\alpha$. Значит, $\alpha$ -- алгебраический.
    
    \item $M = \{g(x) \in F[x]\;|\; g(\alpha) = 0\}$. $M$ непусто, т.к. $\alpha$ -- алгебраический.\\
    
    Докажем, что $M$ -- простой идеал в $F[x]$:
    
    \begin{enumerate}
        \item $g, h \in M \Rightarrow (g + h)(\alpha) = g(\alpha) + h(\alpha) = 0 + 0 = 0$
        \item $g \in M, h \in F[x] \Rightarrow (gh)(\alpha) = g(\alpha)h(\alpha) = 0$ (т.к. $g(\alpha) = 0$)
        \item $(gh) \in M \Rightarrow g(\alpha)h(\alpha) = (gh)(\alpha) = 0 \Rightarrow g(\alpha) = 0$ или $h(\alpha) = 0$, т.е. один из них $\in M$
    \end{enumerate}
    
    Тогда $F(\alpha) \cong F[\alpha] \cong F[x]/(m_{\alpha}(x))$
\end{itemize}

\EndProof